\section{Conclusion et perspectives}
\label{sec:conclusion}

\subsection{Bilan technique}

\TODO{Synthétiser :
\begin{itemize}
    \item Les exigences ET1-ET9 validées (ou partiellement validées, à détailler honnêtement)
    \item Les performances finales : taux de reconnaissance, temps de traitement, empreinte mémoire
    \item Le respect de la contrainte Arduino Due obligatoire
\end{itemize}}

\subsection{Apports pédagogiques}

\TODO{Décrire ce que le projet a apporté :
\begin{itemize}
    \item Mise en pratique de l'ADC, ISR, timer sur SAM3X8E
    \item Conception de filtres numériques temps réel
    \item Pipeline MFCC complet en C
    \item Conception, training et export d'un CNN
    \item Inférence embarquée sans framework lourd
    \item Gestion de projet en équipe, méthodologie cycle en V
\end{itemize}}

\subsection{Perspectives d'évolution}

\TODO{Proposer des évolutions :
\begin{itemize}
    \item \textbf{Élargissement du vocabulaire} : passer de 2 mots (« bleu », « rouge ») à $\geq$~10 mots avec un modèle plus large
    \item \textbf{Affichage OLED} : remplacer les LEDs par un écran SSD1306 (I²C) pour afficher le mot reconnu et la confiance
    \item \textbf{Mode streaming} : détecter automatiquement les mots sans appui bouton (wake-word detection)
    \item \textbf{Bluetooth} : transfert des MFCC vers un smartphone pour traitement cloud
    \item \textbf{Quantification int8} : réduire l'empreinte du modèle et accélérer l'inférence
    \item \textbf{Multi-langue} : étendre le dataset à d'autres langues
    \item \textbf{Robustesse accrue} : augmentation de données plus agressive (réverbération, pitch-shift)
\end{itemize}}

\subsection{Mot de la fin}

\TODO{Paragraphe court de clôture : satisfaction générale, remerciements éventuels (encadrants, camarades).}

\newpage
